\documentclass[11pt,a4paper]{article}

% --- packages ---
\usepackage[utf8]{inputenc}
\usepackage[T2A]{fontenc}
\usepackage[russian]{babel}
\usepackage{lmodern}
\usepackage[a4paper,margin=2.5cm]{geometry}
\usepackage{amsmath,amssymb}
\usepackage{amsthm}
\usepackage[final,expansion=false]{microtype}
% allow texttt to break on slashes/hyphens to avoid right-margin overflow
\usepackage[htt]{hyphenat}
\sloppy
\usepackage{booktabs}
\usepackage{array}
\usepackage{tabularx}
\usepackage{enumitem}
\usepackage[hidelinks]{hyperref}
\usepackage{xcolor}
\usepackage{listings}

% --- listings setup (ASCII-only) ---
\lstset{
  basicstyle=\ttfamily\small,
  breaklines=true,
  columns=fullflexible,
  showstringspaces=false,
  upquote=true,
  literate={'}{{\textquotesingle}}1
}

% --- amsthm style (russian theorem names) ---
\theoremstyle{plain}
\newtheorem{proposition}{Предложение}
\newtheorem{theorem}{Теорема}
\newtheorem{corollary}{Следствие}

% --- metadata ---
\title{GoldenFloat: семейство форматов с плавающей точкой со статическим разбиением,\\
выведенным из золотого сечения \(\varphi\), от GF4 до GF256\\
с точным целочисленным тождеством Люка}
\author{Дмитрий Васильев\\
Независимый исследователь, Ко Самуи, Таиланд\\
Эл. почта: admin@t27.ai \quad ORCID: 0009-0008-4294-6159}
\date{2026-06-02 \quad v1.9}

\hypersetup{
  pdftitle={GoldenFloat: семейство форматов с плавающей точкой со статическим разбиением, выведенным из золотого сечения},
  pdfauthor={Dmitrii Vasilev},
  pdfsubject={Семейство числовых форматов; статическое разбиение, выведенное из phi; тождество Люка},
  pdfkeywords={GoldenFloat, phi, золотое сечение, числа Люка, плавающая точка, posit, takum, OCP-MX}
}

\begin{document}

% ======================================================================
% Журнальная рубрика / УДК (для журнала «Программирование», ИСП РАН)
% ======================================================================
\begin{flushleft}\small
Рубрика: Языки, компиляторы, системы программирования и аппаратно-зависимые вычисления\\[2pt]
УДК 004.222
\end{flushleft}

\maketitle

% Даты заполняются редакцией:
\begin{flushleft}\small
Поступила в редакцию: \rule{2.5cm}{0.4pt}\quad
После доработки: \rule{2.5cm}{0.4pt}\quad
Принята к публикации: \rule{2.5cm}{0.4pt}\\[2pt]
\textbf{DOI:} \rule{4cm}{0.4pt} \emph{(присваивается редакцией и издательством)}
\end{flushleft}

% ======================================================================
% ABSTRACT
% ======================================================================
\begin{abstract}
\noindent\textbf{Что сообщается в этой статье.}
Мы представляем аппаратно-ориентированное описание GoldenFloat~(GF) ---
семейства форматов с плавающей точкой со статическим разбиением, порождаемого
единым замкнутым правилом, --- и три конкретных артефакта: (i)~открытый
многоразрядный RTL-генератор, охватывающий GF4--GF256, с дифференциальной
прогонкой в режиме непрерывной интеграции относительно корректно округлённого
эталона; (ii)~целочисленный тракт аккумулятора, точный по числам Люка, проверенный с
точностью до 500 знаков для \(n = 1,\dots,256\); и (iii)~кодек GF16 на FPGA,
проходящий тестовое окружение 35 из 35 на частоте 323~МГц на Artix-7
(Xilinx XC7A35T). Оракул соответствия формата (Corona) поставляется в том же
репозитории и используется как проверка чёрного ящика в нашем аудите
непрерывной интеграции.

\noindent\textbf{Правило и область его применения.}
Для каждой полной разрядности \(N \ge 4\) ширина поля порядка равна
\(e = \mathrm{round}((N-1)/\varphi^{2})\) при ширине поля мантиссы
\(f = N-1-e\) и \(\varphi = (1+\sqrt{5})/2\). Правило воспроизводит
реализованные ширины поля порядка девяти форматов GF4, GF8, GF12, GF16, GF20,
GF24, GF32, GF64, GF256 (9/9) и согласованно распространяется на GF6, GF10,
GF14, GF48, GF96, GF128, GF512, GF1024. Правило позиционируется рядом с
posit (стандарт Posit 2022 года),
takum (Hunhold 2024, 2025), OCP-MX (Rouhani и др.\ 2023) и черновиком
многоразрядного формата с плавающей точкой IEEE P3109 --- все они являются
семействами, охватывающими ряд разрядностей под параметризованным правилом.
Мы не делаем никаких заявлений о точности или превосходстве на отдельной
ступени по сравнению с любым из них.

\noindent\textbf{Что остаётся открытым.}
Тезис о широте/согласованности инструментальной цепочки зафиксирован как
открытая гипотеза с предрегистрированным путём фальсификации: эксперимент на
FPGA на согласованном субстрате и программная абляция при согласованном
бюджете. Реестр фальсификации (FL-002) фиксирует открытые вопросы и
эксперименты, которые их разрешили бы. Errata о корректности RTL от 2026-05-31
сообщается в разделе~\ref{sec:hw-erratum}; изготовленные кристаллы TTSKY26b
несут дефектный портфель умножителей, а исправленный генератор является
базовой версией для перегенерации.
\end{abstract}

\noindent\textbf{Ключевые слова:} числовые форматы; плавающая точка; золотое
сечение; числа Люка; posit; takum; OCP-MX; арифметическое аппаратное
обеспечение.

% ======================================================================
% 1. INTRODUCTION
% ======================================================================
\section{Введение}

Числовые форматы для машинного обучения и научных вычислений
диверсифицировались далеко за пределы IEEE~754. Доминирующая ось проектирования ---
это то, как фиксированный бюджет из \(N\) бит разбивается между динамическим
диапазоном (поле порядка) и точностью (поле мантиссы) и как это разбиение
масштабируется при изменении \(N\). Три точки сравнения задают рамки настоящей
заметки.

\paragraph{Posit.}
Gustafson и Yonemoto (2017) ввели posit с кодированием «режим плюс порядок»,
параметризованным размером поля порядка \emph{es}. В литературе встречаются
два различных расписания \emph{es}: достандартное расписание с
\(es = 0,1,2,3,4\) при разрядностях \(8,16,32,64,128\) (de Dinechin и др., 2019)
и ратифицированный стандарт Posit 2022 года, который фиксирует \(es = 2\) для
всех разрядностей. Таким образом, динамический диапазон posit масштабируется
вместе с полем режима, которое само является функцией представляемого
значения.

\paragraph{Takum.}
Hunhold (2024, 2025) усовершенствовал идею конусного (tapered) формата при
помощи единого вогнутого правила проекции, определённого единообразно для всех
длин слова. Проекция takum отображает полный целочисленный диапазон \(N\)-битного
слова на интервал числовой прямой; ширина поля мантиссы варьируется
поэлементно, а динамический диапазон масштабируется с \(N\) в рамках одного
аналитического отображения. Takum обратно совместим с IEEE-754 при
\(N \in \{16,32,64\}\) в том смысле, что оболочка динамического диапазона
содержит соответствующий двоичный тип IEEE, и он несёт единственную известную
нам прошедшую рецензирование многоразрядную аппаратную реализацию (ARITH 2025).

\paragraph{OCP-MX.}
Rouhani и др. (2023) стандартизировали блочно-масштабируемые 4-, 6- и 8-битные
типы микромасштабирования. Здесь поэлементная мантисса коротка, но общий
масштаб уровня блока несёт динамический диапазон. Семейство микромасштабирования
является развёрнутым промышленным эталоном для форматов малой разрядности.

\paragraph{Место GoldenFloat.}
GoldenFloat исследует намеренно иную точку в этом пространстве проектирования:
\emph{статическое разбиение}, фиксированное для каждой разрядности и не
зависящее от хранимого значения. Поэтому динамический диапазон GF\(N\)
определяется полностью шириной поля порядка \(e(N)\) и выбором смещения, а не
каким-либо зависящим от значения кодированием.
Это обменивает поэлементную адаптивность posit и takum на свойство, которое
труднее получить в значение-адаптивных семействах: замкнутое аналитическое
правило, единообразно фиксирующее разбиение порядок/мантисса по всей лестнице
разрядностей, плюс целочисленное тождество аккумулятора, зависящее только от
выбора основания, а не от конкретной разрядности.

\paragraph{Историческая преемственность.}
Работа продолжает отечественную линию исследований оригинальных систем
представления чисел в вычислительной технике, начатую троичной ЭВМ
«Сетунь» под руководством Н.~П.~Брусенцова (Вычислительный центр МГУ,
1958), --- первой в мире машиной на симметричном троичном коде
\((-1, 0, +1)\), а в «Сетунь-70» (1970) --- с введённым троичным форматом
«трайт» и оптимизированным диапазоном мантиссы нормализованного числа.
GoldenFloat переносит этот дух поиска неканонических числовых
представлений в современный контекст форматов малой разрядности для
машинного обучения. По доступным открытым источникам, GoldenFloat --- один
из первых именованных числовых форматов российского авторства, доведённых
до уровня публикации и аппаратной реализации (FPGA).

Мотивация разбиения бюджета в точке \(1/\varphi^{2}\) опирается на три факта,
ни один из которых не является оригинальным для данной работы:

\begin{itemize}\setlength{\itemsep}{2pt}
\item Алгебраическое соотношение \(\varphi^{2} = \varphi + 1\) устраняет
постоянный множитель в кодере на золотом сечении (Golden Ratio Encoder)
Daubechies, Gunturk, Wang и Yilmaz (IEEE TIT, 2010): бета-кодер с основанием
\(\varphi\) допускает безмножительную реализацию робастного АЦП только на
сложениях. Это инженерная причина предпочесть основание \(\varphi\) внутри
диапазона робастных оснований; это \emph{не} утверждение о том, что \(\varphi\)
является единственным допустимым основанием.
\item То же алгебраическое соотношение даёт классическое тождество Люка
\(\varphi^{2n} + \varphi^{-2n} = L_{2n}\) (Lucas, 1878; Бине),
которое позволяет накапливать \(\varphi\)-масштабированные частичные суммы через
целочисленную рекуррентность Люка. Это арифметическая причина, по которой
\(\varphi\)-основанная лестница допускает целочисленный тракт аккумулятора без
аппаратного полутипа.
\item Отношение \(1/\varphi^{2} \approx 0.382\) попадает внутрь эмпирического
интервала отношений, который воспроизводит реализованные ширины поля порядка GF
(раздел~\ref{sec:ladder}). Это отношение разделяют 82 других значения в данном
интервале; выбор именно \(\varphi\) мотивирован двумя приведёнными выше
алгебраическими фактами, а не одним лишь численным воспроизведением.
\end{itemize}

Вклад данной заметки состоит из трёх аппаратно-программных артефактов плюс
замкнутого лестничного правила, связывающего их воедино:

\begin{itemize}\setlength{\itemsep}{2pt}
\item \textbf{(C1) Открытый многоразрядный RTL-генератор.} Единый
параметризованный генератор на Verilog выпускает кодеки и умножители для всей
лестницы \(N \ge 4\) (GF4--GF256) из одного шаблона \((E, M, \mathrm{BIAS})\),
с шириной регистра произведения \([2M+1:0]\) и округлением полупозиции вверх
(round-half-up).
Исправленный генератор (раздел~\ref{sec:hw-erratum}) является базовой версией
для перегенерации и встроен в дифференциальную прогонку непрерывной интеграции
относительно корректно округлённого эталона; см. раздел~\ref{sec:hw-status}.
\item \textbf{(C2) Целочисленный тракт аккумулятора, точный по числам Люка.}
Классическое тождество Люка \(\varphi^{2n} + \varphi^{-2n} = L_{2n}\)
(Lucas 1878; Бине), проверенное символически и с точностью до 500 десятичных
знаков для \(n = 1,\dots,256\), позволяет нести \(\varphi\)-масштабированные
частичные суммы в беззнаковом целочисленном хранилище без аппаратного полутипа.
Мы не претендуем на это тождество как на своё; мы заявляем инженерное
следствие: оно допускает целочисленный аккумулятор без quire в стиле posit
или регистра Кулиша (Kulisch) (раздел~\ref{sec:lucas}).
\item \textbf{(C3) Оракул соответствия формата (Corona), поставляемый как
встроенный в дерево чёрный ящик.} Corona --- это детерминированный эталон для
каждого формата, который проверяет каждый кодек и умножитель, выпускаемый (C1),
относительно замкнутой спецификации (C2). Он используется как шлюз CI в
репозитории GoldenFloat и является единственной точкой, обнаружившей дефект
портфеля умножителей от 2026-05-31 (раздел~\ref{sec:hw-corona},
раздел~\ref{sec:hw-erratum}).
\end{itemize}

Лестничное правило \(e = \mathrm{round}((N-1)/\varphi^{2})\), \(f = N-1-e\)
(\(N \ge 4\)) точно воспроизводит девять реализованных ширин поля порядка GF
(9/9) и согласованно распространяется на GF6, GF10, GF14, GF48, GF96, GF128,
GF512, GF1024. Битстрим GF16 на FPGA
проходит тестовое окружение кодека 35 из 35 на частоте 323~МГц на Artix-7
(Xilinx XC7A35T); ни один физический кристалл не вернулся ни для одной ступени.

Мы рассматриваем это как исследовательскую заметку, а не как программный
документ. Мы не утверждаем, что какая-либо ступень GF точнее равнонастроенного
не-\(\varphi\) формата той же разрядности; не утверждаем, что \(\varphi\) ---
единственное или привилегированное основание; и не заявляем о кремниевой
валидации на какой-либо ступени. Там, где свидетельства неполны, мы говорим об
этом, помечаем как открытое и приводим эксперимент, который это
фальсифицировал бы. Раздел~\ref{sec:ladder} определяет лестничное правило и
сообщает об исчерпывающем поиске с поправкой на множественность (look-elsewhere).
Раздел~\ref{sec:gre} фиксирует якорь кодера на золотом сечении и точно
указывает, что он устанавливает, а что нет.
Раздел~\ref{sec:lucas} проверяет тождество Люка и мотивирует целочисленный
аккумулятор.
Раздел~\ref{sec:hw} сообщает о статусе аппаратной части по лестнице GF4--GF256.
Раздел~\ref{sec:prior} соотносит GF с предшествующими работами.
Раздел~\ref{sec:prereg} даёт предрегистрацию открытых гипотез.
Раздел~\ref{sec:ledger} даёт реестр фальсификации FL-002.
Приложения содержат верификационный скрипт, индекс форматов, таблицу поправки
на множественность, предрегистрацию FPGA-эксперимента на согласованном
субстрате и протокол репликатора.

% ======================================================================
% 2. LADDER RULE
% ======================================================================
\section{Лестничное правило \texorpdfstring{$e = \mathrm{round}((N-1)/\varphi^{2})$}{e = round((N-1)/phi^2)}}\label{sec:ladder}

\subsection{Определение}
Пусть \(N \ge 4\) --- полная разрядность формата в битах, включая один знаковый
бит. Разбиение GoldenFloat выделяет
\begin{equation}
e = \mathrm{round}\!\left(\frac{N-1}{\varphi^{2}}\right), \qquad
f = N - 1 - e,
\label{eq:ladder}
\end{equation}
где \(\varphi = (1+\sqrt{5})/2 = 1.6180339887\ldots\) и
\(\varphi^{2} = \varphi + 1 = 2.6180339887\ldots\),
так что \(1/\varphi^{2} = 0.3819660\ldots\).
Разбиение \emph{статично}: \(e\) зависит только от \(N\) и никогда от
представляемого значения. В таблице~\ref{tab:ladder} приведён результат для
семнадцати форматов (девять реализованных разрядностей плюс восемь
разрядностей-расширений). При \(N \to \infty\)
\(\mathrm{round}((N-1)/\varphi^{2})/(N-1) \to 1/\varphi^{2}\), и реализованное
отношение в таблице~\ref{tab:ladder} сходится соответственно. Для \(N \ge 4\)
нет граничной неустойчивости.

\paragraph{Краевые случаи при \(N = 2,3\).}
Применение правила при \(N = 2\) даёт \(e = 0\) --- знаковое однобитное поле без
порядка, что не является представлением с плавающей точкой в обычном смысле.
При \(N = 3\) оно даёт \(e = 1\), \(f = 1\) --- четырёхзначный формат с
пренебрежимо малыми диапазоном и точностью. Поэтому мы определяем лестницу как
от GF4 до GF256 и рассматриваем GF2/GF3 как вырожденные краевые случаи формулы,
а не как реализованные форматы. Все утверждения настоящей заметки касаются
\(N \ge 4\).

\begin{table}[t]
\centering
\caption{Лестничное правило GF для семнадцати форматов. Столбцы: полная
разрядность \(N\), биты порядка \(e\), биты мантиссы \(f = N-1-e\), сырое
значение \((N-1)/\varphi^{2}\) до округления и реализованное отношение
\(e/(N-1)\) (цель \(1/\varphi^{2} = 0.38197\)). Верхний блок перечисляет девять
разрядностей с вернувшимся кремнием или финализированным RTL. Средний блок
(GF6, GF10, GF14, GF48, GF96, GF128) перечисляет выведенные из правила ступени
без вернувшегося кремния; RTL для GF128 написан, но не отправлен на
изготовление. Нижний блок (GF512, GF1024) сообщает о двух ступенях-расширениях
правила, чьи смещения (\(2^{194}-1\), \(2^{390}-1\)) превышают поле
\texttt{u128}, используемое в записи ROM Corona, и отслеживаются символически
только на уровне оракула t27 SSOT. Все семнадцать строк арифметически проверены
по уравнению~\eqref{eq:ladder}. Значения вычислены с точностью 200 знаков в
\texttt{mpmath}.}
\label{tab:ladder}
\begin{tabular}{rrrrr}
\toprule
$N$ & $e$ & $f$ & $(N-1)/\varphi^{2}$ & $e/(N-1)$ \\
\midrule
4    & 1   & 2   & 1.1459   & 0.33333 \\
8    & 3   & 4   & 2.6738   & 0.42857 \\
12   & 4   & 7   & 4.2016   & 0.36364 \\
16   & 6   & 9   & 5.7295   & 0.40000 \\
20   & 7   & 12  & 7.2574   & 0.36842 \\
24   & 9   & 14  & 8.7852   & 0.39130 \\
32   & 12  & 19  & 11.8409  & 0.38710 \\
64   & 24  & 39  & 24.0639  & 0.38095 \\
256  & 97  & 158 & 97.4013  & 0.38039 \\
\midrule
6    & 2   & 3   & 1.9098   & 0.40000 \\
10   & 3   & 6   & 3.4377   & 0.33333 \\
14   & 5   & 8   & 4.9656   & 0.38462 \\
48   & 18  & 29  & 17.9524  & 0.38298 \\
96   & 36  & 59  & 36.2868  & 0.37895 \\
128  & 49  & 78  & 48.5097  & 0.38583 \\
\midrule
512  & 195 & 316 & 195.1846 & 0.38160 \\
1024 & 391 & 632 & 390.7512 & 0.38221 \\
\bottomrule
\end{tabular}
\end{table}

\subsection{Поправка на множественность (look-elsewhere) для совпадения с \texorpdfstring{$\varphi^{-2}$}{phi^-2}}
\label{sec:lookelsewhere}

\paragraph{Сначала отрицательный результат.}
\(\varphi^{-2} \approx 0.38197\) невозможно отличить от прочих 82 совпадающих
отношений на основании представленных здесь численных свидетельств. Приведённая
ниже поправка на множественность (look-elsewhere) делает это точным.

\paragraph{Постановка.}
Мы искали отношения \(r \in [0.1, 0.9]\) с шагом \(10^{-5}\), что даёт мощность
пространства поиска \(N_s = 80\,000\). Каждое отношение проверялось против
\(W = 9\) реализованных разрядностей GoldenFloat; отношение \emph{совпадает},
если оно воспроизводит все 9 разрядностей. Поиск вернул \(K = 83\) совпадающих
отношения.

\paragraph{Семейственная вероятность.}
При нулевой гипотезе о том, что каждое отношение проходит тест каждой разрядности
независимо с равномерной поразрядной вероятностью \(p\), откалиброванной так,
что \(\mathbb{E}[\text{совпадений}] = K\), имеем
\(p_{\text{match}} = K/N_s = 83/80\,000 \approx 1.04 \times 10^{-3}\).
Число совпадений \(X \sim \mathrm{Binomial}(N_s, p_{\text{match}})\).
Вычисляя через регуляризованную неполную бета-функцию с точностью 60 знаков:
\[
P(X \ge 83) \approx 7.1 \times 10^{-3}.
\]
Таким образом, семейственная вероятность наблюдения по меньшей мере 83
совпадений при данной нулевой гипотезе не пренебрежимо мала; совпадающее
множество является умеренно частым исходом поиска, а не ярким событием в хвосте
распределения.

\paragraph{Поправка Бонферрони для \(\varphi^{-2}\).}
Нескорректированная вероятность того, что любое заранее заданное отношение
попадёт в совпадающее множество, равна
\(p_{\text{match}} \approx 1.04 \times 10^{-3}\).
После поправки Бонферрони для \(N_s = 80\,000\) одновременных тестов
скорректированное p-значение насыщается на единице (поскольку
\(N_s \cdot p_{\text{match}} = K = 83\)).
\(\varphi^{-2}\) не значимо по Бонферрони.

\paragraph{Сужение с 12 форматами.}
Расширение поиска на все двенадцать форматов в таблице~\ref{tab:ladder} сужает
множество кандидатов с 392 отношений (девятиформатный интервал
\([0.37844, 0.38235]\)) до 47 отношений (двенадцатиформатный интервал
\([0.38189, 0.38235]\)) --- сокращение в \(8.3\times\). Единственность при более
тонком разрешении требует точной интервальной арифметики и здесь не заявляется.

\paragraph{Вывод.}
Численное совпадение того, что \(\varphi^{-2}\) удовлетворяет лестничному
правилу, реально, но не доказательно: 82 других отношения удовлетворяют тому же
правилу. Аргументация в пользу \(\varphi\) опирается на его структурные свойства
(алгебраическое самоподобие, рекуррентность Фибоначчи, точное по Люка
накопление), а не на одно лишь совпадение в пространстве поиска. Все вычисления
воспроизводимы через \texttt{look\_elsewhere\_calc.py} в дополнительных
материалах.

\subsection{Режим округления}
Правило использует округление к ближайшему чётному, но этот выбор несущественен
для каждой реализованной разрядности, как фиксирует сноска.\footnote{Для всех
девяти реализованных разрядностей сырое значение \((N-1)/\varphi^{2}\) не
является полуцелым (ближайшее приближение --- \(N = 64\) при \(24.064\) и
\(N = 128\) при \(48.510\)), поэтому выбор режима округления не влияет ни на одну
реализованную разрядность. Тем не менее мы задаём округление к ближайшему чётному
(по умолчанию в IEEE 754) для формальной полноты.}

% ======================================================================
% 3. GRE ANCHOR
% ======================================================================
\section{Безмножительный якорь GRE}\label{sec:gre}

Независимая, прошедшая рецензирование мотивация для \(\varphi\) исходит из
кодера на золотом сечении (Golden Ratio Encoder, GRE) Daubechies, Gunturk, Wang
и Yilmaz, ``The Golden Ratio Encoder'', \emph{IEEE Transactions on Information
Theory} 56(10):5097--5110, 2010 (arXiv:0809.1257). GRE --- это схема
аналого-цифрового преобразования, а не формат с плавающей точкой. То, что он
устанавливает и что мы заимствуем, --- это точное инженерное свойство основания
\(\varphi\):

\begin{quote}
\(\varphi\) --- это основание, при котором бета-кодер допускает безмножительную
реализацию робастного АЦП только на сложениях, поскольку
\(\varphi^{2} = \varphi + 1\) устраняет компоненту с постоянным множителем
(Daubechies и др., IEEE TIT 2010).
\end{quote}

\subsection{Что результат GRE устанавливает, а что нет}

Мы явно очерчиваем область действия, поскольку два естественных перетолкования
ложны.

\paragraph{Нет утверждения, что выбор основания вынужден.}
Мы не утверждаем, что \(\varphi\) --- единственное допустимое основание,
дающее работоспособный кодер. В той же статье анализируются другие значения
\(\beta\) и демонстрируется диапазон устойчивых и робастных оснований. Любое
утверждение о том, что \(\varphi\) является единственным допустимым основанием,
было бы ложным и здесь не делается.

\paragraph{Нет точечной робастности при \(\beta = \varphi\).}
Чистый алгоритмический кодер при \(\beta = \varphi\) устойчив, но не робастен.
Робастность установлена только для реализации с диапазоном параметров /
ориентированным ациклическим графом (Daubechies и др., теорема 8), а не для
единственной точки \(\beta = \varphi\). Мы не смешиваем это с анализом
восстановления \(\beta\) у Ward (2008, arXiv:0806.1083), который является
другой статьёй, обращающейся к другому вопросу, и не является источником
безмножительного свойства, используемого здесь.

Поэтому связь, которую мы проводим, узка и защитима:
\(\varphi^{2} = \varphi + 1\) --- это алгебраический факт, устраняющий
постоянный множитель в рекурсии GRE, и тот же алгебраический факт лежит в основе
тождества Люка из раздела~\ref{sec:lucas}. Это инженерная причина предпочесть
\(\varphi^{2}\) внутри совпадающего интервала из раздела~\ref{sec:lookelsewhere},
а не утверждение о том, что кодер однозначно выбирает \(\varphi\).

% ======================================================================
% 4. LUCAS IDENTITY
% ======================================================================
\section{Точное целочисленное тождество Люка}\label{sec:lucas}

\subsection{Мотивация: большие накопления и целочисленные суммы}
\label{sec:lucas-motivation}

Практическая забота для числовых форматов малой разрядности --- поведение
потери значимости при больших накопленных суммах. В арифметическом тракте с
фиксированной точностью с плавающей точкой и шириной мантиссы \(f\)
относительная погрешность длинного тракта умножения с накоплением растёт по мере
увеличения числа накопленных слагаемых: каждое сложение округляет текущую сумму
до \(f+1\) значащих бит, а абсолютная погрешность масштабируется (в худшем
случае) с величиной частичной суммы. Это хорошо понятая мотивация за
конструкциями точного накопления, такими как \emph{quire} в posit (Gustafson и
Yonemoto, 2017) и аккумулятор Кулиша (Kulisch), оба из которых дополняют
плавающий тракт внешним широким регистром с фиксированной точкой, чтобы длинные
суммы могли вычисляться без промежуточного округления.

Арифметическое свойство, которое мы используем в
разделе~\ref{sec:lucas-statement} ниже, даёт иной путь к той же цели,
специфичный для основания \(\varphi\): \(\varphi\)-масштабированные частичные
суммы вида \(\sum_i \varphi^{2 n_i}\) можно отслеживать через \emph{целочисленную}
рекуррентность по числам Люка, причём сопряжённый остаток
\(\sum_i \varphi^{-2 n_i}\) точен и ограничен. Аппаратное следствие в том, что
длинный аккумулятор, построенный вокруг основания \(\varphi\), не требует
внешнего широкого регистра с фиксированной точкой или аппаратного полутипа;
целочисленное хранилище и рекуррентность Люка сами являются аккумулятором.

Мы излагаем это здесь как инженерную причину интереса к тождеству \emph{до}
доказательства и возвращаемся в разделе~\ref{sec:lucas-engineering} к тому, что
заявляется и что не заявляется о конкретном аппаратном аккумуляторе.

\subsection{Формулировка и классическое происхождение}\label{sec:lucas-statement}

\begin{proposition}[Точное по Люка тождество для $\varphi^{2}$]
\label{prop:lucas}
Для каждого целого \(n \ge 1\),
\begin{equation}
\varphi^{2n} + \varphi^{-2n} = L_{2n},
\label{eq:lucas}
\end{equation}
где \(L_k\) --- \(k\)-е число Люка (\(L_0 = 2\), \(L_1 = 1\),
\(L_k = L_{k-1} + L_{k-2}\)).
\end{proposition}

Это классическое следствие формулы Бине
\(L_k = \varphi^{k} + (-\varphi)^{-k}\): для чётного \(k = 2n\) знаковый множитель
\((-1)^{k}\) обращается, оставляя ровно целое \(L_{2n}\). Случай \(n = 1\) даёт
якорь
\begin{equation}
\varphi^{2} + \varphi^{-2} = 3 = L_{2},
\label{eq:trinity-anchor}
\end{equation}
который мы приписываем Lucas (1878) и формуле Бине и никогда не представляем как
оригинальный для данной работы или её автора.

\subsection{Верификация (F1)}\label{sec:lucas-verification}

Мы проверили уравнение~\eqref{eq:lucas} двумя независимыми способами для
\(n = 1,\dots,256\) (так что \(2n = 2,\dots,512\)):

\begin{enumerate}\setlength{\itemsep}{2pt}
\item \textbf{Символически.} Используя точную арифметику в
\(\mathbb{Q}[\sqrt{5}]\) (SymPy), алгебраическое тождество
\(\varphi^{2n} + \varphi^{-2n} - L_{2n} = 0\) выполняется для каждого \(n\) в
диапазоне.
\item \textbf{Численно.} При 500 десятичных знаках (\texttt{mpmath})
максимальный остаток \(|\varphi^{2n} + \varphi^{-2n} - L_{2n}|\) по диапазону
равен \(1.55 \times 10^{-499}\) при \(n = 256\), значительно ниже целочисленного
масштаба. Это уровень численного шума, согласующийся с точностью 500 знаков.
\end{enumerate}

Все 256 точных проверок SymPy проходят тождественно; наихудший относительный
остаток равен \(1.55 \times 10^{-499}\) при \(n = 256\). Сокращённая таблица
результатов (избранные строки от \(n = 1\) до \(n = 256\)) приведена в
приложении~A. Скрипт воспроизведён в приложении~A.

\subsection{Инженерное следствие и что не заявляется}
\label{sec:lucas-engineering}

Предложение~\ref{prop:lucas} означает, что степени \(\varphi^{2}\) удовлетворяют
линейной целочисленной рекуррентности. Следовательно, \(\varphi\)-масштабированные
частичные суммы вида \(\sum_i \varphi^{2 n_i}\) в принципе можно накапливать в
беззнаковом целочисленном хранилище, отслеживая соответствующие целые числа Люка
\(L_{2 n_i}\), при этом остаточные члены \(\varphi^{-2 n_i}\) отслеживаются
символически или ограничиваются. Это арифметическая основа целочисленного тракта
накопления, не зависящего от аппаратного полутипа и не требующего внешнего
широкого регистра с фиксированной точкой, такого как quire или аккумулятор
Кулиша.

Мы ясно заявляем, что не утверждается. Конкретный аппаратный аккумулятор,
реализующий эту идею, с заданными длинами слов, границами переполнения,
задержкой и сравнением с quire posit или аккумулятором Кулиша, является будущей
работой и здесь не заявляется. Само тождество проверено
(раздел~\ref{sec:lucas-verification}); инженерный артефакт --- нет.
Релевантная предшествующая работа по целочисленной арифметике в базисе
Фибоначчи/\(\varphi\) --- это Ahlbach, Usatine и Pippenger, ``Efficient
Algorithms for Zeckendorf Arithmetic'' (2012, arXiv:1207.4497), которую мы
цитируем как якорь алгоритмической осуществимости целочисленной арифметики в
базисе \(\varphi\).

% ======================================================================
% 5. HARDWARE STATUS
% ======================================================================
\section{Статус аппаратной части по лестнице}\label{sec:hw}

\subsection{Определение лестницы готовности}
Мы сообщаем о статусе аппаратной части относительно шестиступенчатой лестницы
готовности, взятой дословно из документа STATUS репозитория
\texttt{tt-trinity-phi}:
\begin{enumerate}\setlength{\itemsep}{1pt}
\item SPEC --- существует версионированная числовая/поведенческая спецификация.
\item RTL --- синтезируемый Verilog, реализующий спецификацию, зафиксирован в репозитории.
\item SIM --- функциональная симуляция (cocotb + iverilog) проходит канонический якорь.
\item SYNTH --- синтез Yosys + чистый линт Verilator.
\item GDS/TAPEOUT --- поток OpenLane2 (SKY130A) производит артефакт tt\_submission; отправка на shuttle принята.
\item SILICON --- физический кристалл измерен и охарактеризован на плате.
\end{enumerate}
Более высокая ступень заявляется только если все нижние ступени удовлетворены с
доказательствами внутри публичного репозитория.

\subsection{Статус по разрядностям}\label{sec:hw-status}
Таблица~\ref{tab:hwstatus} сообщает о статусе по разрядностям, проверенном
прямой инспекцией ветки master \texttt{gHashTag/t27} и ветки main
\texttt{gHashTag/tt-trinity-phi} 2026-05-31.
В таблице ``TTSKY26b'' обозначает принятие в отправку на shuttle TTSKY26b
(SKY130A); счётчики байт для каждого RTL-файла опущены ради вёрстки и
зафиксированы в репозитории.

\begin{table}[t]
\centering
\caption{Статус аппаратной части по разрядностям по лестнице GF, проверенный
инспекцией репозитория (\texttt{gHashTag/t27} master и
\texttt{gHashTag/tt-trinity-phi} main, 2026-05-31).
Строки, помеченные \(\dagger\), имели дефект умножителя на уровне генератора в
RTL, отправленном в TTSKY26b 2026-05-17; дефект раскрыт в
разделе~\ref{sec:hw-erratum}, а исправленный RTL опубликован в
\texttt{gHashTag/tt-trinity-gamma} PR\#110 и \texttt{gHashTag/t27} PR\#1024.}
\label{tab:hwstatus}
\footnotesize
\setlength{\tabcolsep}{4pt}
\begin{tabularx}{\linewidth}{@{}l >{\raggedright\arraybackslash}X l l l l@{}}
\toprule
Разрядность & RTL-файл(ы) & SIM & SYNTH & GDS shuttle & Кремний \\
\midrule
GF4   & \texttt{gf4\_add.v}   & только RTL & только RTL & TTSKY26b & не вернулся \\
GF8   & \texttt{gf8\_add.v}$^\dagger$   & только RTL & только RTL & TTSKY26b & не вернулся \\
GF12  & \texttt{gf12\_add.v}$^\dagger$  & только RTL & только RTL & TTSKY26b & не вернулся \\
GF16  & \texttt{gf16\_add.v}, \texttt{gf16\_dot4.v}, \texttt{gf16\_mul.v}$^\dagger$; битстрим FPGA, Artix-7, 35/35 на 323~МГц & да (0x47C0) & да & TTSKY26b$^\ast$ & не вернулся \\
GF20  & \texttt{gf20\_add.v}$^\dagger$  & только RTL & только RTL & TTSKY26b & не вернулся \\
GF24  & \texttt{gf24\_add.v}$^\dagger$  & только RTL & только RTL & TTSKY26b & не вернулся \\
GF32  & \texttt{gf32\_add.v}$^\dagger$  & только RTL & только RTL & TTSKY26b & не вернулся \\
GF64  & \texttt{gf64\_add.v}$^\dagger$  & только RTL & только RTL & TTSKY26b & не вернулся \\
GF128 & \texttt{gf128\_add.v}$^\dagger$ & только RTL & только RTL & TTSKY26b & не вернулся \\
GF256 & \texttt{gf256\_add.v}$^\dagger$ & только RTL & только RTL & TTSKY26b & не вернулся \\
\bottomrule
\multicolumn{6}{@{}l@{}}{\scriptsize $^\ast$\,GF16 был также отправлен ранее в shuttle TTSKY26a.} \\
\end{tabularx}
\end{table}

Разрядность GF16 --- единственная ступень с документированным битстримом FPGA и
функциональной симуляцией, проходящей канонический якорь скалярного произведения
0x47C0. Все десять разрядностей от GF4 до GF256 несут синтезируемый RTL на
Verilog-2005 в \texttt{gHashTag/tt-trinity-phi/src/} и были включены в отправку
на shuttle TTSKY26b в Tiny Tapeout (SKY130A, отправлено 2026-05-17). Shuttle
TTSKY26a (ядро сетки dot4 GF16) был отправлен ранее с изготовлением,
запланированным на 2026-10-28. Ни один физический кристалл не был возвращён,
охарактеризован или измерен ни для одной ступени; поэтому любая цифра «на ватт»,
«на джоуль» или пиковой пропускной способности вне архитектурных проекций
является спроектированной, а не измеренной.
DOI аппаратного архива \texttt{10.5281/zenodo.19227877} цитируется только как
архив RTL и артефактов отправки на shuttle, но никогда как ссылка на
производительность или результаты.

\subsection{Оракул соответствия формата (Corona)}\label{sec:hw-corona}

Сопутствующий кристалл только для чтения, \texttt{gHashTag/tt-trinity-corona},
готовится для shuttle TTGF26a Tiny Tapeout (GF180MCU 180~нм, 4\(\times\)4
тайла, целевая отправка 2026-06-22, ожидаемый кремний 2026-10 -- 2026-11).
Corona --- четвёртый кристалл в линейке TRI-NET после Phi, Euler и Gamma.
Это \emph{оракул соответствия формата}, а не вычислительный ускоритель: запрос
приходит как 7-битный индекс формата на \texttt{ui\_in[6:0]}, и чип возвращает
запрошенные поля записи внутрикристального ROM, кодирующего все 80 записей
числовых форматов единого источника истины TRI-NET
(\texttt{gHashTag/t27 specs/numeric/formats\_catalog.t27}, PR~\#1028),
вместе с 17 уникальными модулями RTL-декодеров Tier-1, обслуживающими 22
внутрикристальных индекса форматов (пять индексов разделяют декодеры, например
FP8~E4M3 с MXFP8~E4M3 и NF4-BNB с NF4-QLoRA), каждый из которых преобразует
внутрикристальный формат в IEEE 754 FP32 или INT32. 80 записей разбиты на
тринадцать кластеров форматов (двоичный IEEE, десятичный IEEE, ML
низкой точности, GoldenFloat, posit/unum-III, OCP-MX, LNS, целый/фиксированный,
исторический, теоретический, сжатие, расширенный, квантово-настроенный). Текущая
конструкция синтезируется приблизительно в 2{,}308 ячеек (\(\approx 1\%\)
бюджета сайтов 4\(\times\)4). HEAD главной ветки репозитория на момент данной
ревизии --- \texttt{b5462ac49892}.

Мы заявляем, чем Corona \emph{не является}, чтобы сохранить честность настоящей
заметки. Corona не делает никаких заявлений о вычислительной пропускной
способности, энергии на операцию или производительности относительно любого
другого чипа-ускорителя, и её существование \textbf{не} является
свидетельством того, что \(\varphi\)-лестница или любой формат GoldenFloat
превосходит конкурирующую числовую систему. Тезис о широте-как-рве
(breadth-as-moat) FL-002 из раздела~\ref{sec:ledger} остаётся [Открытая
гипотеза], и Corona не меняет его статус. Takum (Hunhold 2024,
\texttt{arXiv:2412.20273}) остаётся действующим живым контрпримером к FL-002 и
поставляется в ROM Corona как запись Tier-2 (или Tier-1, если лицензирование
VHDL эталонной реализации Hunhold разрешится благоприятно); он не подавляется.
Чип по построению только для чтения, поэтому никакой класс дефекта округления
вида, раскрытого в разделе~\ref{sec:hw-erratum}, не может применяться к его
основному тракту чтения из ROM; 17 декодеров Tier-1 подвергаются тому же
дифференциальному прогону RTL-аудита, что и остальной портфель, с пятью слоями
доказательства на декодер (исчерпывающая прогонка, независимый эталон, формальная
эквивалентность, постсиликоновый вектор, kill-мутации), покрывающими в общей
сложности 592{,}308 входных кодов. На момент настоящей рукописи (50
направленных тестов PASS, GDS и предпроверка Tiny~Tapeout PASS, 21~CI-шлюз
автообнаружен) готовность Corona находится на ступени GDS/TAPEOUT лестницы
раздела~\ref{sec:hw} в ожидании окна отправки TTGF26a; ни один кремний Corona не
был измерен.

\subsection{Конфаундер аппаратного субстрата для программных сравнений BPB}
\label{sec:hw-confound}

Сравнения бит-на-байт (bits-per-byte, BPB) между GF16 и IEEE fp16,
сообщаемые конвейером обучения IGLA RACE
(\texttt{gHashTag/trios-trainer-igla}), производятся на субстратах ---
потребительских GPU и трактах CPU SIMD, --- чьё аппаратное умножение с
накоплением спроектировано и настроено под IEEE fp16~/~bf16, с нативными
TensorCore или AMX, диспетчеризующими fp16/bf16 за один такт. GF16 на том же
субстрате эмулируется через целочисленно-кодированные тракты без нативного
умножителя. Поэтому сообщаемый разрыв BPB порядка \(10^{-1}\) между GF16 и
fp16 на этом субстрате не измеряет формат GF16 --- он измеряет формат GF16 на
аппаратуре, оптимизированной против него. Справедливое сравнение BPB требует
либо fp16-нативного кремния, спаренного с GF16-нативным кремнием при
согласованном технологическом узле и площади, либо контролируемой FPGA-оснастки,
где оба формата получают эквивалентный бюджет LUT/DSP. Пока такое спаренное
измерение не существует, программно-субстратные дельты BPB в любом направлении
фиксируются как искажённые субстратом и не являются свидетельством за или против
широты-как-рва. Этот открытый вопрос зафиксирован как пункт~(a) FL-002 в
разделе~\ref{sec:ledger} и является основной причиной, по которой ни одно
заявление о точности на отдельной ступени нигде в данной заметке не появляется.

\subsection{Errata о корректности RTL}\label{sec:hw-erratum}

2026-05-31 дифференциальная прогонка на уровне RTL относительно корректно
округлённого эталона с плавающей точкой обнаружила систематический дефект на
уровне генератора в большинстве умножителей GoldenFloat в виде, отправленном на
shuttle TTSKY26b Tiny Tapeout 2026-05-17 (полная помодульная разбивка, вывод
корневой причины и обсуждение объёма исправления приведены в
приложении~\ref{app:erratum}). Регистр произведения в отправленном RTL
умножителя был объявлен на два бита уже, поэтому для базового входа
\(1.0 \times 1.0\) старший бит был усечён, и результат читался как \(0.5\),
реплицированный по всему сгенерированному портфелю. Исправленный универсальный
генератор умножителя с шириной произведения \([2M+1:0]\) проходит исчерпывающие
прогонки на gf8 (26{,}360/26{,}360), gf12 (109{,}576/109{,}576), gf16
(262{,}144/262{,}144), gf20 (100{,}000/100{,}000), gf24
(100{,}000/100{,}000), gf32 (200{,}000/200{,}000) и направленные точные тесты
на gf64/gf128/gf256 и теперь встроен в CI-шлюз
(\texttt{run\_gf\_audit.sh}) при каждом коммите. Кремний TTSKY26b уже находится
в изготовлении и не может быть отозван, поэтому кристаллы gamma и phi будут
нести дефектный RTL умножителя; исправленный портфель является базовой версией
для перегенерации для следующего shuttle, и ни одно заявление о точности на
отдельной ступени в данной заметке не зависело от умножителей в виде, в каком они
были отправлены (раздел~\ref{sec:hw-status}). Дефект не отменяет лестничную
арифметику раздела~\ref{sec:ladder}, тождество Люка раздела~\ref{sec:lucas} или
учёт множественности раздела~\ref{sec:ladder}, все из которых являются
абстрактными спецификациями, не зависящими от какой-либо конкретной реализации
RTL.

\subsection{Эмпирический вердикт на реальном корпусе}\label{sec:hw-realcorpus}

Программное прямое сопоставление между конвейером смешанной точности на
\(\varphi\)-лестнице и гетерогенным «зоопарком» числовых форматов было выполнено
на корпусе реального текста (\texttt{tiny\_shakespeare} --- наименьший публично
аудируемый корпус, который мы смогли зафиксировать) внутри оснастки IGLA RACE
2026-05-31. Сначала и прямо фиксируем агрегированный вердикт: сравнение даёт
\emph{недостаточно свидетельств}. Данные не подтверждают заявление о
превосходстве на отдельной ступени ни в одном из направлений.

\paragraph{Агрегированный результат (бит-на-байт; меньше --- лучше).}
На \texttt{tiny\_shakespeare} среднее BPB на отложенной выборке составило
\(\mathrm{BPB}_{\varphi} = 5.9871\) (выборочное стандартное отклонение
\(\sigma = 0.0911\)) для плеча \(\varphi\)-лестницы и
\(\mathrm{BPB}_{\mathrm{zoo}} = 6.0454\) (\(\sigma = 0.2083\)) для плеча
гетерогенного зоопарка по \(n < 11\) парным сидам. Двухрешенческий вердиктный
комплект был таков: (i)~сравнение фронта Парето при согласованном бюджете бит на
вес возвращает \emph{зоопарк побеждает} (точка зоопарка \(P_w = 8.0\) достигает
\(\mathrm{BPB} = 5.361\), \(\mathrm{CI} = [5.348, 5.375]\), тогда как точка
\(\varphi\)-лестницы \(P_w = 4.0\) достигает \(\mathrm{BPB} = 5.360\),
\(\mathrm{CI} = [5.347, 5.373]\), причём два доверительных интервала
перекрываются); (ii)~нормализованный байесовский вторичный критерий возвращает
ничью; (iii)~третичная апостериорная достоверность возвращает
\(P(\varphi < \mathrm{zoo}) = 0.976\), ниже предрегистрированного порога
решения для заявления о точности на отдельной ступени. Поэтому агрегированный
вердикт разрешается в недостаточно свидетельств. Мы не утверждаем, что
\(\varphi\)-лестница точнее зоопарка при согласованном бюджете бит; мы также не
утверждаем обратного.

\paragraph{Открытые пробелы, от которых зависит вердикт.}
Три пробела не позволяют повысить вердикт с «недостаточно свидетельств» до
проверенного результата в любую сторону, каждый зафиксирован в реестре
фальсификации (раздел~\ref{sec:ledger}, подпункт FL-002):
(1)~плечо зоопарка не диспетчеризует по \(P_w\);
(2)~нормализованная байесовская достоверность вычисляется, но не встроена в
агрегат; и
(3)~регрессионная фиксация (regression pin) отсутствует. Размер выборки в этом
прогоне (\(n < 11\) парных сидов) ниже целевого минимально обнаружимого эффекта
\(n = 11\), требуемого для \(\mathrm{MDE} = 0.05\) BPB при \(\alpha = 0.05\),
\(\text{мощность} = 0.80\).

% ======================================================================
% 6. PRIOR ART
% ======================================================================
\section{Связь с предшествующими работами}\label{sec:prior}

GoldenFloat находится среди нескольких опубликованных усилий охватить диапазон
разрядностей единым параметризованным правилом. Мы рассматриваем их как
предшественников и союзников, а не как конкурентов, которых следует отвергнуть,
и не утверждаем, что GF лучше любого из них по какой-либо оси.

\paragraph{Posit.}
Gustafson и Yonemoto (2017) ввели posit с кодированием «режим плюс порядок»,
параметризованным размером поля порядка \emph{es}. В литературе встречаются два
различных расписания \emph{es}: достандартное расписание с
\(es = 0,1,2,3,4\) при разрядностях \(8,16,32,64,128\) (de Dinechin и др., 2019)
и ратифицированный стандарт Posit 2022 года, который фиксирует \(es = 2\) для
всех разрядностей. Аппаратные реализации включают PERCIVAL (Mallasen и др.,
2022) и Big-PERCIVAL (Mallasen и др., 2024). Лестница posit является
естественным контролем для вопроса широты (раздел~\ref{sec:ledger}).

\subsection{GoldenFloat против Takum: структурированное сравнение бок о бок}
\label{sec:gf-takum}

\paragraph{Признание приоритета.}
Арифметика takum (Hunhold 2024, 2025) --- ближайшая опубликованная альтернатива
лестнице GoldenFloat с прошедшей рецензирование многоразрядной аппаратной
реализацией. Статья Hunhold на ARITH 2025 представляет RTL- и FPGA-результаты для
нескольких разрядностей takum под единым замкнутым правилом проекции, опережая
любое сопоставимое аппаратное заявление для GoldenFloat. Лестница GoldenFloat в
настоящее время не имеет бенчмарка точности на отдельной ступени против takum, не
имеет прошедшей рецензирование аппаратной публикации и не имеет прямого
численного сравнения. Таблица~\ref{tab:gftakum} фиксирует, что делает каждый
формат и где стоит каждое заявление; это не конкурентный рейтинг.

В настоящее время: takum имеет прошедшие рецензирование аппаратные результаты на
FPGA на нескольких разрядностях; GoldenFloat --- нет. GoldenFloat заявляет
свойство целочисленного аккумулятора, точного по Люка; takum это не нацелено.
Ни один из форматов не имеет прямого бенчмарка численной точности против другого.

\begin{table}[t]
\centering
\caption{Лестница GoldenFloat против takum (Hunhold 2024) --- структурированное
сравнение. Это не конкурентный рейтинг; превосходство по точности не заявляется.}
\label{tab:gftakum}
\small
\begin{tabularx}{\linewidth}{l X X}
\toprule
Ось & Лестница GoldenFloat (GF) & Takum (Hunhold 2024) \\
\midrule
Тип разбиения & Статическое: ширины полей порядка и мантиссы фиксированы для каждой ступени замкнутым правилом \(e = \mathrm{round}((N-1)/\varphi^{2})\); не зависят от значения. & Конусное (tapered): единая вогнутая проекция отображает полный целочисленный диапазон на интервал числовой прямой; ширина мантиссы варьируется поэлементно. \\
Единое замкнутое правило & Да --- лестничное правило \(N \mapsto (e, f)\) детерминировано; нет поэлементного ветвления на уровне формата. & Да --- проекция takum --- это единое аналитическое отображение, применяемое единообразно к каждому \(n\)-битному целому. \\
Базовая константа & Золотое сечение \(\varphi = (1+\sqrt{5})/2\); ширины поля порядка растут как дроби последовательных чисел Фибоначчи от полной разрядности слова. & Основание натурального логарифма \(e\): takum --- логарифмический формат конусной точности, чьё кодированное значение является функцией \(\ln(x)\), с полями знак + режим + характеристика + мантисса (Hunhold 2024). \\
Статус аппаратной части & RTL-портфель существует (исправлен после обнаружения дефекта умножителя в дифференциальной прогонке; см. раздел~\ref{sec:hw-erratum}); синтез FPGA сообщён внутренне; нет прошедшего рецензирование кремниевого результата. & RTL- и FPGA-результаты прошли рецензирование на ARITH 2025 (Hunhold 2024/2025); кремний не сообщён. \\
Прошедшая рецензирование многоразрядная аппаратная часть & Отсутствует на момент написания. Это документированный открытый пробел. & Да --- Hunhold ARITH 2025 охватывает несколько разрядностей под единым правилом с числами синтеза FPGA. Это действующий эталон. \\
Целочисленно-точный аккумулятор & GF сообщает о точном по Люка тождестве аккумулятора: \(\varphi^{2n} + \varphi^{-2n} = L_{2n}\) (\(n = 1,\dots,256\) при 500 знаках; раздел~\ref{sec:lucas}). Это свойство специфично для выбора константы \(\varphi\). & Аналогичное целочисленно-точное тождество аккумулятора не документировано. Отсутствие не является дефектом; оно отражает иную цель проектирования. \\
\bottomrule
\end{tabularx}
\end{table}

\paragraph{IEEE P3109.}
Рабочая группа IEEE P3109 с 2023 года разрабатывает черновик стандарта
многоразрядного двоичного формата с плавающей точкой, нацеленного на ускорители
машинного обучения. Текущий публичный рабочий черновик (v0.9.1, 2025)
специфицирует двоичные форматы с плавающей точкой на тринадцати разрядностях (от
3 до 15 бит) под единым параметризованным правилом, с эталонной реализацией,
поддерживаемой Graphcore (\texttt{graphcore-research/gfloat}), и формальной
верификацией, сообщённой на ARITH 2025. P3109 --- ближайшее
\emph{действующее, отраслевое} семейство форматов с плавающей точкой,
охватывающее ряд разрядностей, к лестнице GoldenFloat и является прямой целью
фальсификации для любого прочтения настоящей работы в духе широты-как-рва
(раздел~\ref{sec:ledger}, FL-002~(a)).

\paragraph{OCP-MX.}
Rouhani и др. (2023) определили форматы микромасштабирования Open Compute
Project --- отраслевое стандартное семейство блочно-масштабируемых 4-, 6- и
8-битных типов. Это развёрнутая эталонная точка для семейств малой разрядности.

\paragraph{LNS.}
Логарифмические системы счисления --- логарифмический предшественник для систем
не по основанию 2, с недавними экземплярами для машинного обучения в LNS-Madam
(2021) и Alam, Garland и Gregg (2021). LNS даёт точное умножение ценой
приближённого сложения; оно не нацелено на целочисленно-точное тождество
аккумулятора в смысле раздела~\ref{sec:lucas}.

\paragraph{Арифметика Цекендорфа.}
Bergman (1957) ввёл систему счисления с \(\varphi\) в качестве иррационального
основания --- исторический предшественник числовых систем на основе \(\varphi\).
Ahlbach, Usatine и Pippenger (2012) дают эффективные алгоритмы для арифметики в
базисе Фибоначчи/\(\varphi\) --- алгоритмическую предшествующую работу для идеи
целочисленного накопления раздела~\ref{sec:lucas}.
Более недавнее предложение с учётом \(\varphi\) --- квантование Fibbinary
(Belghazi и др., 2025, arXiv:2511.01921), которое представляет веса в базисе
Фибоначчи/\(\varphi\) и сообщает о серьёзной деградации точности без
восстановления через обучение с учётом квантования. Fibbinary работает на уровне
поэлементного кодирования веса, а не как полное семейство форматов с плавающей
точкой, поэтому оно дополняет GoldenFloat, а не является прямой альтернативой.

\paragraph{Аккумулятор Кулиша.}
Широкий аккумулятор с фиксированной точкой Кулиша (Kulisch, 2013) и его
воплощение в виде \emph{quire} posit (стандарт Posit 2022) --- канонический
эталон для точного накопления скалярного произведения. Точный по Люка тракт
GoldenFloat не является аккумулятором Кулиша: он остаётся внутри базиса
\(\varphi\) и использует тождество \(\varphi^{2n} + \varphi^{-2n} = L_{2n}\) для
сохранения частичных сумм в беззнаковом целочисленном хранилище, тогда как Кулиш
хранит явный широкий регистр с фиксированной точкой и преобразует в него и из него.

Для более широкого контекста обзор числовых систем DNN (Alsuhli и др., 2023),
определения форматов FP8 (Micikevicius и др., 2022), Huawei HiFloat8
(Luo и др., 2024), числовые форматы в климатических моделях (Kloewer и др., 2020),
анализ масштабных законов точности (Kumar и др., 2024) и монография Gustafson
\emph{The End of Error} (2015) помещают эти семейства в более широкий ландшафт.
На масштабе менее 1 млрд параметров BitNet b1.58 (Ma и др., 2024) и
последующая работа 2B4T (Ma и др., 2025, arXiv:2504.12285) являются текущими
чемпионами по бит-на-байт и естественным базисом ниже 1 млрд, относительно
которого любое заявление о рецепте обучения GoldenFloat должно в конечном счёте
измеряться.

\paragraph{Честное резюме.}
Нам не известно о другом опубликованном семействе форматов с плавающей точкой,
которое выводит своё разбиение \(e:f\) из единого замкнутого правила по лестнице
от GF4 до GF256 с целочисленным тождеством, точным по Люка. Posit, takum,
OCP-MX, LNS и черновик многоразрядного формата с плавающей точкой IEEE P3109 ---
близкие предшественники и союзники, охватывающие ряд разрядностей собственными
едиными параметризованными правилами; takum несёт ближайшее прошедшее
рецензирование многоразрядное аппаратное заявление, а IEEE P3109 несёт ближайшее
усилие отраслевой стандартизации на малых разрядностях. Специфичный для
GoldenFloat элемент --- это комбинация статического разбиения \(e:f\) из одного
замкнутого правила и целочисленного тракта аккумулятора, точного по Люка, который
зависит только от выбора основания \(\varphi\), а не от разрядности.

% ======================================================================
% 7. PREREGISTRATION
% ======================================================================
\section{Предрегистрация открытых гипотез}\label{sec:prereg}

\paragraph{Дата и область авторства.}
Эта предрегистрация была составлена 2026-05-31 до того, как появились
какие-либо кремниевые измерения при согласованном бюджете. Она заморожена на
git-теге \texttt{v1.9-prereg} в \texttt{gHashTag/goldenfloat-preprint}
(см. запись CHANGELOG для v1.9 для SHA-1 коммита) и не будет изменена после
начала сбора данных при согласованном бюджете.

Четыре зарегистрированные гипотезы и их условия остановки таковы.

\paragraph{$H_a$ (Широта-как-ров выдерживает контроль при согласованном бюджете).}
Нулевая: дельта BPB на бит на площадь кристалла между портфелем GoldenFloat и
fp16-нативным базисом равного транзисторного бюджета равна нулю или отрицательна
для GoldenFloat. Условие остановки: дельта BPB при согласованном бюджете падает
ниже \(\delta_{\text{thresh}} = 0.005\) бит-на-байт на согласованном корпусе
оценки; объявляется фальсификация, и тезис широты-как-рва в
разделе~\ref{sec:hw-confound} отзывается.

\paragraph{$H_b$ (Лестничное правило единственно внутри пространства поиска).}
Нулевая: по меньшей мере одно альтернативное правило отбора, построенное без
обращения к лестнице GoldenFloat, воспроизводит те же 12 канонических
разрядностей. Условие остановки: интервал поиска для сертификата единственности
сужается до синглтонного множества; любое нетривиальное конкурирующее правило,
опубликованное до v1.5, вносится в реестр фальсификации.

\paragraph{$H_c$ (Смещение GF256 допускает вывод в замкнутой форме).}
Нулевая: не существует выражения в замкнутой форме, воспроизводящего константу
смещения GF256 с полной двойной точностью IEEE-754 в аудируемой символической
форме. Условие остановки: точная формула смещения помещена в аудируемой форме
(вывод CAS + независимая выборочная проверка) в дополнительных материалах;
формула должна пройти верификацию SymPy с точностью 500 знаков.

\paragraph{$H_d$ (Кремний с исправленным RTL GF16 достигает паритета измерений).}
Нулевая: исправленный кристалл GF16 и fp16-нативный эталонный кристалл,
согласованные по технологическому узлу и бюджету площади, различаются более чем
на \(X = 2\%\) по погрешности накопления туда-обратно (round-trip) на
согласованных тестовых векторах. Условие остановки: парные измерения кристаллов
из независимого испытательного дома попадают в пределы \(X\) процентов; пока не
существует вернувшегося кристалла, эта гипотеза остаётся открытой, и нигде в
препринте не делается кремниевого заявления.

% ======================================================================
% 8. LEDGER FL-002
% ======================================================================
\section{Открытые вопросы и реестр фальсификации FL-002}\label{sec:ledger}

Мы фиксируем открытые вопросы как реестр фальсификации, указывая для каждого
статус, заявление и эксперимент, который его фальсифицировал бы. Отрицательные
результаты излагаются первыми и прямо.

\paragraph{(a) Широта-как-ров (открытая гипотеза).}
Гипотеза состоит в том, что преимущество лестницы GF --- это широта и
согласованность инструментальной цепочки (одно замкнутое правило по всей
лестнице), а не точность на отдельной ступени. Это не доказано абляцией, и
единственные доступные программные данные BPB искажены субстратом
(раздел~\ref{sec:hw-confound}). Фальсификация:
(i)~справедливое спаренное измерение --- GF16-нативный кремний против
fp16-нативного кремния при согласованном технологическом узле и площади, или
FPGA-оснастка с эквивалентными бюджетами LUT/DSP для обоих форматов --- в котором
лестница posit, takum, MX или fp16 соответствует лестнице GF при согласованных
бюджетах бит;
(ii)~F2, прямое сопоставление бит-на-байт конвейера смешанной точности
\(\varphi\)-лестницы против гетерогенного зоопарка при согласованных бюджетах
бит, с учётом потерь при межформатных преобразованиях;
(iii)~F3, контроль лестницы posit при тех же бюджетах бит.
Любой отрицательный исход по (i)--(iii) должен быть зафиксирован прямо и первым.

\paragraph{(b) Неединственность лестничного правила (риск).}
Как показано в разделе~\ref{sec:lookelsewhere}, 83 значения отношений
воспроизводят девять проверенных разрядностей; совпадающий интервал ---
\([0.3786, 0.3822]\), и \(1/\varphi^{2}\) внутри него не выделяется.
Фальсификация (или сужение): расширить множество проверенных разрядностей
(например, добавить GF128 и GF512) и сообщить, сужается ли совпадающий интервал
до синглтона или остаётся множественным.

\paragraph{(c) Хранимое смещение GF256 (открыто).}
Хранимое смещение приблизительно \(2^{71}\) документировано для GF256 против
ожидания в стиле IEEE \(2^{96}-1\) для 97-битного поля порядка. Это
расхождение не объяснено в аудируемой форме. Фальсификация (или разрешение):
опубликовать точную формулу смещения и её вывод.

\paragraph{(d) Кремний за пределами GF16 (открыто).}
Как документирует раздел~\ref{sec:hw-status}, все десять разрядностей GF4 --
GF256 несут синтезируемый RTL и были включены в отправку на shuttle TTSKY26b в
Tiny Tapeout, но ни один физический кристалл не вернулся ни для одной ступени.
Фальсификация (или разрешение) --- возврат кремния и функциональная проверка на
отдельной ступени против оракула соответствия формата
раздела~\ref{sec:hw-corona}.

\paragraph{(e) Портфель умножителей в виде, отправленном в TTSKY26b (отозван).}
Как документируют раздел~\ref{sec:hw-erratum} и приложение~\ref{app:erratum},
RTL умножителя, отправленный на shuttle TTSKY26b 2026-05-17, несёт единственную
ошибку формулы генератора (регистр произведения на два бита уже),
реплицированную по разрядностям, помеченным крестом (\(\dagger\)). Поэтому
изготовленные кристаллы gamma и phi будут нести дефектный портфель умножителей.
Исправленный универсальный генератор является базовой версией для перегенерации;
статус отозван для отправки TTSKY26b, открыт в ожидании возврата кремния для
следующего shuttle.

\paragraph{(f) FPGA-эксперимент на согласованном субстрате $H_4$ (запланирован).}
Кодеки GF16 и posit16, синтезированные на одном и том же Xilinx Artix-7
(XC7A100T, QMTech Wukong V1) с согласованными бюджетами LUT/DSP, будут оценены
по числу LUT, числу триггеров, \(F_{\max}\) и динамической мощности на операцию
умножения с накоплением. Предрегистрированные пороги
(приложение~\ref{app:fpga}): \(R_{\mathrm{LUT}} < 1.15\) и
\(R_{F_{\max}} > 0.85\). Провал по любому порогу понижает широту-как-ров с
открытой гипотезы до риска по оси кремния. Любой отрицательный результат
сообщается первым и заметно.

% ======================================================================
% 9. CONCLUSION
% ======================================================================
\section{Заключение}\label{sec:conclusion}

Лестничная арифметика воспроизводит девять реализованных ширин поля порядка
9 из 9 и согласованно распространяется на восемь дальнейших выведенных из
правила ступеней --- GF6, GF10, GF14, GF48, GF96, GF128, GF512, GF1024
(раздел~\ref{sec:ladder}); тождество Люка
\(\varphi^{2n} + \varphi^{-2n} = L_{2n}\) выполняется точно для
\(n = 1,\dots,256\), с якорем \(\varphi^{2} + \varphi^{-2} = 3 = L_{2}\),
приписываемым Lucas (1878) и Бине (раздел~\ref{sec:lucas}); и единственный кодек
GF16 проходит тестовое окружение 35 из 35 на частоте 323~МГц на Artix-7
(раздел~\ref{sec:hw-status}). Широта-как-ров --- заявление о том, что одно
правило по всей лестнице даёт преимущество согласованности, а не преимущество
точности на отдельной ступени, --- остаётся открытой гипотезой
(раздел~\ref{sec:ledger}\,(a)). Что её фальсифицировало бы: лестница posit,
takum или MX, соответствующая \(\varphi\)-лестнице при согласованных бюджетах
бит. FPGA-эксперимент на согласованном субстрате (приложение~D,
предрегистрированный) обеспечит первое контролируемое аппаратное сравнение. Два
дальнейших пункта открыты (смещение GF256 и кремний за пределами GF16) и один
является риском (неединственность лестничного правила, с 83 совпадающими
отношениями). Мы представляем эту заметку, чтобы установить цитируемую,
фальсифицируемую запись проверенной арифметики и открытых вопросов.

% ======================================================================
% APPENDICES
% ======================================================================
\appendix

\section{Верификационный скрипт F1 и расширенная таблица}\label{app:f1}

Приведённый ниже скрипт проверяет \(\varphi^{2n} + \varphi^{-2n} = L_{2n}\)
символически (SymPy, точно в \(\mathbb{Q}[\sqrt{5}]\)) и численно
(\texttt{mpmath}, 500 знаков) для \(n = 1,\dots,256\).

\begin{lstlisting}[language=Python]
#!/usr/bin/env python3
"""F1: verify phi^(2n) + phi^(-2n) = L_{2n} (integer Lucas number)
for n = 1..256, both symbolically (sympy, exact) and numerically
(mpmath, 500 dps). Anchor: Ahlbach/Usatine/Pippenger (2012),
arXiv:1207.4497. The identity is a classical Binet corollary
(Lucas, 1878), not original to this work."""
from mpmath import mp, sqrt, power, nstr, mpf
from sympy import sqrt as ssqrt, simplify, Integer

def lucas_recurrence(max_index):
    L = [2, 1]
    for _ in range(2, max_index + 1):
        L.append(L[-1] + L[-2])
    return L[: max_index + 1]

def main():
    mp.dps = 500
    phi = (1 + sqrt(5)) / 2
    n_max = 256                        # gives 2n = 2..512
    L = lucas_recurrence(2 * n_max)
    max_diff = mpf(0)
    all_pass = True
    for n in range(1, n_max + 1):
        m = 2 * n
        val = power(phi, m) + power(phi, -m)
        diff = abs(val - L[m])
        max_diff = max(max_diff, diff)
        all_pass = all_pass and (diff < mpf("1e-150"))
    # exact symbolic check in Q[sqrt(5)]
    phi_sym = (Integer(1) + ssqrt(5)) / 2
    exact_pass = all(
        simplify(phi_sym**(2*n) + phi_sym**(-2*n) - Integer(L[2*n])) == 0
        for n in range(1, n_max + 1))
    print("Anchor phi^2 + phi^-2 =",
          nstr(power(phi, 2) + power(phi, -2), 20), "= L_2 =", L[2])
    print("max |phi^(2n)+phi^(-2n) - L_2n| =", nstr(max_diff, 5))
    print("all numerical rows within 1e-150 ?", all_pass)
    print("exact symbolic identity holds for n=1..256 ?", exact_pass)
    return 0 if all_pass else 1

if __name__ == "__main__":
    raise SystemExit(main())
\end{lstlisting}

\begin{table}[t]
\centering
\caption{Расширенный вывод F1. \(\varphi^{2n} + \varphi^{-2n}\) равно целому
числу Люка \(L_{2n}\) с точностью 500 знаков. Избранные строки из полного
прогона \(n = 1,\dots,256\); остатки абсолютные.}
\label{tab:f1-extended}
\small
\begin{tabular}{rrrr}
\toprule
$n$ & $2n$ & $L_{2n}$ & $|\text{остаток}|$ \\
\midrule
1   & 2   & 3                         & 0 \\
2   & 4   & 7                         & $1.39 \times 10^{-501}$ \\
4   & 8   & 47                        & $1.66 \times 10^{-501}$ \\
8   & 16  & 2207                      & $4.53 \times 10^{-501}$ \\
16  & 32  & 4870847                   & $8.4 \times 10^{-501}$ \\
32  & 64  & 23725150497407            & $1.99 \times 10^{-500}$ \\
64  & 128 & $5.629 \times 10^{26}$    & $3.75 \times 10^{-500}$ \\
128 & 256 & $3.168 \times 10^{53}$    & $7.67 \times 10^{-500}$ \\
192 & 384 & $1.783 \times 10^{80}$    & $1.15 \times 10^{-499}$ \\
256 & 512 & $1.004 \times 10^{107}$   & $1.55 \times 10^{-499}$ \\
\bottomrule
\end{tabular}
\end{table}

\section{Индекс форматов GF (от GF4 до GF256)}\label{app:format-index}

Таблица~\ref{tab:format-index} даёт для каждой реализованной разрядности: полное
число бит \(N\), биты порядка \(e\), биты мантиссы \(f\), отношение \(e/(N-1)\)
(цель \(1/\varphi^{2} = 0.38197\)), phi-расстояние \(|E/M - 1/\varphi|\),
где \(E/M = e/f\) и \(1/\varphi = 0.61803\), и реализована ли разрядность.

\begin{table}[h]
\centering
\caption{Индекс форматов GF для девяти реализованных разрядностей.}
\label{tab:format-index}
\begin{tabular}{rrrrrc}
\toprule
$N$ & $e$ & $f$ & $e/(N-1)$ & $|E/M - 1/\varphi|$ & реализ.? \\
\midrule
4   & 1  & 2   & 0.33333 & 0.11803 & Да \\
8   & 3  & 4   & 0.42857 & 0.13197 & Да \\
12  & 4  & 7   & 0.36364 & 0.04661 & Да \\
16  & 6  & 9   & 0.40000 & 0.04863 & Да \\
20  & 7  & 12  & 0.36842 & 0.03470 & Да \\
24  & 9  & 14  & 0.39130 & 0.02482 & Да \\
32  & 12 & 19  & 0.38710 & 0.01354 & Да \\
64  & 24 & 39  & 0.38095 & 0.00265 & Да \\
256 & 97 & 158 & 0.38039 & 0.00411 & Да \\
\bottomrule
\end{tabular}
\end{table}

\section{Таблица поправки на множественность}\label{app:lookelsewhere}

Таблица~\ref{tab:lookelsewhere} сообщает для репрезентативных правил-кандидатов,
сколько из девяти реализованных разрядностей
(\(N \in \{4,8,12,16,20,24,32,64,256\}\), с целевыми значениями \(e\)
\(\{1,3,4,6,7,9,12,24,97\}\)) каждое воспроизводит. Исчерпывающий рациональный
поиск по \(p/q\) с \(p \in \{1,\dots,99\}\), \(q \in \{100,\dots,499\}\) нашёл
83 различных значения отношений, совпадающих со всеми девятью; совпадающий
интервал --- \([0.3786, 0.3822]\). Правило GF
\(\mathrm{round}((N-1)/\varphi^{2})\) --- один член этого множества и не
выделяется одним лишь воспроизведением (раздел~\ref{sec:lookelsewhere}).

\begin{table}[h]
\centering
\caption{Правила-кандидаты и их счёт воспроизведения девяти разрядностей.}
\label{tab:lookelsewhere}
\begin{tabular}{lll}
\toprule
Правило & Совпадения & Примечание \\
\midrule
$\mathrm{round}((N-1)/\varphi^{2})$ & 9/9 & правило GF; отношение $\approx 0.38197$ \\
$\lfloor N/\varphi^{2} \rfloor$     & 9/9 & столь же просто; то же отношение \\
$\mathrm{round}((N-1) \cdot 0.382)$ & 9/9 & округлённая константа; без $\varphi$ \\
$\mathrm{round}((N-1) \cdot 3/7.85)$ & 9/9 & произвольная рациональная; без $\varphi$ \\
$\mathrm{round}((N-1) \cdot 3/8)$    & 8/9 & не проходит GF256 (96 против 97) \\
$\mathrm{round}((N-1) \cdot 5/13)$   & 8/9 & не проходит GF256 \\
$\lfloor N \cdot 3/8 \rfloor$        & 8/9 & не проходит GF256 \\
$\mathrm{round}((N-1)/2.6)$          & 8/9 & не проходит GF256 \\
$\mathrm{round}((N-1)/e)$            & 5/9 & не проходит $N$=24,32,64,256 \\
$\lfloor (N-1)/\varphi^{2} \rfloor$  & 5/9 & floor против round имеет значение \\
$\mathrm{round}((N-1)/\pi)$          & 2/9 & плохое соответствие \\
$\mathrm{round}((N-1)/\varphi)$      & 0/9 & неверная константа \\
\bottomrule
\end{tabular}
\end{table}

\section{FPGA-эксперимент на согласованном субстрате (предрегистрированный)}\label{app:fpga}

\paragraph{Гипотеза $H_4$.}
Кодеки GF16 и posit16, синтезированные на одном и том же Xilinx Artix-7
(XC7A100T-FGG676, QMTech Wukong V1) с согласованными бюджетами LUT/DSP,
достигают сопоставимых числа LUT, числа триггеров, максимальной рабочей частоты
(\(F_{\max}\)) и динамической мощности на операцию умножения с накоплением.

\paragraph{Предрегистрированные пороги.}
Пусть \(R_{\mathrm{LUT}} = \mathrm{LUT}_{\mathrm{GF16}} / \mathrm{LUT}_{\mathrm{posit16}}\)
и \(R_{F_{\max}} = F_{\max,\mathrm{GF16}} / F_{\max,\mathrm{posit16}}\).
Эксперимент фальсифицирует \(H_4\) тогда и только тогда, когда
\[
R_{\mathrm{LUT}} > 1.15 \quad \text{ИЛИ} \quad R_{F_{\max}} < 0.85.
\]

\paragraph{Правило решения.}
Если условие фальсификации выполнено, широта-как-ров понижается с открытой
гипотезы до риска по оси кремния в пункте~(a) FL-002. Отрицательный результат
сообщается первым, в аннотации или резюме отчёта о результатах, без задержки на
опровержение.
Если оба отношения в пределах границ, \(H_4\) сохраняется, и статус
широты-как-рва остаётся открытой гипотезой в ожидании измерения IGLA RACE на
втором субстрате.
Если GF16 превосходит posit16 по площади или частоте, это наблюдение
фиксируется, но не повышает широту-как-ров до проверенной. Один кодек на одном
субстрате --- это не ров.

\paragraph{Материалы.}
RTL GF16: исправленные \texttt{gf16\_mul.v} и \texttt{gf16\_add.v} из
\texttt{gHashTag/tt-trinity-gamma} PR \#110.
RTL posit16: PERCIVAL (Mallasen и др., 2022, arXiv:2111.15286).
Программирование: Rust-драйвер \texttt{cli/dlc10} (\texttt{gHashTag/t27}).
Инструментальная цепочка: OpenXC7 (yosys + nextpnr-xilinx) или Vivado 2023.2.
Шлюз корректности: исчерпывающая прогонка TestFloat-3, требуется 0 ошибок в
1 млн случайных выборок до принятия результатов по площади/таймингу.

\paragraph{Обязательство.}
Любой отрицательный результат --- то есть любой исход, где GF16 требует более
чем на 15\% дополнительной площади или опускается ниже 85\% от \(F_{\max}\)
posit16, --- будет сообщён прямо и заметно. Никакой постфактумной корректировки
порога, никакого избирательного опущения.

Дата предрегистрации: 2026-05-31.
Полная спецификация SHA-256: см. git-тег \texttt{v1.9-prereg} в
\texttt{gHashTag/goldenfloat-preprint}; SHA-256 рукописи на этом теге является
авторитетным хешем предрегистрации.

\section{Протокол независимого репликатора}\label{app:replicator}

Это приложение даёт точные шаги воспроизведения для каждого проверяемого
заявления в препринте. Каждый шаг перечисляет команду, ожидаемый результат и
заполнитель для дайджеста SHA-256.

\paragraph{Шаг R1: Проверка предложения~\ref{prop:lucas} (тождество Люка).}
\texttt{python f1\_lucas\_proof\_extended.py}.
Ожидаемый результат: код выхода 0; вывод терминала включает строку
\texttt{MAX RESIDUAL < 1e-150 at 500 digits} и строку
\texttt{F1 VERIFIED}.

\paragraph{Шаг R2: Проверка лестничного правила (12/12 разрядностей).}
\texttt{python ladder\_extended.py}.
Ожидаемый результат: код выхода 0; вывод терминала включает
\texttt{LADDER 12/12 WIDTHS REPRODUCED} и нет токенов \texttt{FAIL}.

\paragraph{Шаг R3: Проверка битстрима FPGA на QMTech Wukong V1.}
\begin{lstlisting}
git clone https://github.com/gHashTag/t27
cd t27/fpga/vsa/
dlc10 idcode
\end{lstlisting}
Ожидаемый idcode: \texttt{0x13631093}.
\begin{lstlisting}
dlc10 sram gf16_heartbeat_top.bit
\end{lstlisting}
Ожидаемый результат: консоль UART печатает \texttt{GF16 HEARTBEAT OK} в течение
5 секунд после загрузки битстрима.

\paragraph{Шаг R4: Проверка аудита исправленного RTL умножителя.}
\begin{lstlisting}
git clone --branch 110 \
    https://github.com/gHashTag/tt-trinity-gamma tt-trinity-gamma-110
cd tt-trinity-gamma-110
bash verification/run_gf_audit.sh
\end{lstlisting}
Ожидаемый результат: код выхода 0; вывод терминала содержит точную строку
\texttt{GF AUDIT ALL PASS} на последней непустой строке.

% ======================================================================
% APPENDIX F: RTL ERRATUM EXTENDED DETAIL
% ======================================================================
\section{Errata о корректности RTL: расширенные подробности}\label{app:erratum}

Это приложение даёт полную помодульную разбивку, вывод корневой причины и объём
исправления дефекта портфеля умножителей от 2026-05-31, кратко изложенного в
разделе~\ref{sec:hw-erratum}. Дефект был найден дифференциальной прогонкой на
уровне RTL относительно корректно округлённого эталона с плавающей точкой в
репозиториях чипов TRI-NET
(\texttt{gHashTag/tt-trinity-gamma} и \texttt{gHashTag/tt-trinity-phi}),
применённой к сгенерированному портфелю умножителей, отправленному на shuttle
TTSKY26b Tiny Tapeout 2026-05-17.

\paragraph{Затронутые модули.}
\texttt{gf8\_mul} не проходит на \(\approx 95\%\) исчерпывающей прогонки из
26{,}360 точек входа; \texttt{gf12\_mul} не проходит на \(\approx 99\%\)
прогонки из 109{,}576 точек; \texttt{gf20\_mul}, \texttt{gf24\_mul} и
\texttt{gf32\_mul} каждый не проходят на \(100\%\) направленных входов (например,
\(1.0 \times 1.0\) читается как \(0.5\), порядок смещён на единицу);
\texttt{gf16\_mul} корректен только после однострочного расширения регистра
округлённой мантиссы с 9 до 10 бит, и вариант в виде, отправленном, несёт ошибку
переполнения округления; \texttt{gf64\_mul}, \texttt{gf128\_mul} и
\texttt{gf256\_mul} разделяют то же семейство генератора и предположительно
затронуты по той же причине, хотя они ещё не были исчерпывающе прогнаны.
\texttt{gf4\_mul} вырожден (\(e = 1\) не оставляет нормальных порядков) и не
затронут. Сумматоры \texttt{gf8\_add} и \texttt{gf12\_add} имели отдельный
дефект нормализации узкого формата (например, \(0.25 + 0.25\) читается как
\(0\)); более широкие сумматоры \texttt{gf16\_add} и \texttt{gf32\_add} были уже
корректны.

\paragraph{Корневая причина.}
Произведение двух \((M+1)\)-битных мантисс (с неявной ведущей единицей)
имеет ширину \(2M+2\) бита и попадает в \([2^{2M}, 2^{2M+2})\). Отправленный RTL
объявил регистр произведения на два бита уже и выполнял нормализацию на битах,
сдвинутых вниз на два. Для самого базового входа
\(1.0 \times 1.0 = 2^{2M}\) ведущий бит усекается, порядок декрементируется, и
результат читается как \(0.5\). Это единственная ошибка формулы генератора
(позиция старшего бита должна быть \(2M + 1\)), реплицированная по всему
портфелю.

\paragraph{Почему покрытие формальными методами это не поймало.}
База доказательств Coq в репозиториях чипов TRI-NET устанавливает математику
уровня формата GoldenFloat (тождество
\(\varphi^{2}+\varphi^{-2}=3\), тождества кодирования, согласованность
лестничного правила); она не устанавливает ширины бит и округление реализации на
Verilog. Дефект целиком живёт в реализации RTL, в коде вне доказанного ядра. Это
усиливает, а не ослабляет позицию верифицируемости настоящей заметки: машинно
проверенная математика уровня формата необходима, но недостаточна; дифференциальная
прогонка на уровне RTL относительно корректно округлённого эталона должна быть
частью конвейера аудита.

\paragraph{Исправление и текущий статус.}
Исправленный универсальный генератор умножителя выпускает правильный шаблон для
любого \((E, M, \mathrm{BIAS})\): ширина произведения \([2M+1:0]\),
нормализация на \([2M+1]\) или \([2M]\), извлечение мантиссы \([2M:M+1]\) или
\([2M-1:M]\), округление полупозиции вверх с переносом. Исправленный портфель
чист: gf8 исчерпывающая прогонка 0 отказов из 26{,}360; gf12 0 из 109{,}576;
gf16 0 из 262{,}144; gf20 0 из 100{,}000;
gf24 0 из 100{,}000; gf32 0 из 200{,}000; gf64, gf128, gf256 проходят
направленные точные тесты (\(1.0 \times 1.0\), \(1.5 \times 1.5\), \(\dots\)),
хотя полные прогонки остаются будущей работой. Исправленный RTL опубликован как
запросы на слияние \texttt{gHashTag/tt-trinity-gamma}\#110 (54 файла, +2452/-1032)
и \texttt{gHashTag/t27}\#1024 (16 файлов, +558/-10), а методологическая заметка
опубликована как \texttt{gHashTag/trinity-clara}\#20. CI-шлюз
(\texttt{run\_gf\_audit.sh}, встроенный в
\texttt{.github/workflows/gf-audit.yml}) теперь повторно прогоняет
дифференциальную прогонку при каждом коммите.

\paragraph{Объём исправления.}
Кремний TTSKY26b уже отправлен, и изготовление не может быть отозвано; поэтому
изготовленные кристаллы gamma и phi будут нести дефектный RTL умножителя.
Канонический якорь POST 0x47C0 и тракты тождества формата должны быть повторно
проверены на вернувшемся кремнии, но ни одно заявление о точности на отдельной
ступени в данной заметке в любом случае не зависело от портфеля умножителей в
отправленном виде (раздел~\ref{sec:hw-status} документировал все ступени, кроме
GF16, как только-RTL). Исправленный портфель является базовой версией для
перегенерации для следующего shuttle. Дефект не отменяет лестничную арифметику
раздела~\ref{sec:ladder}, тождество Люка раздела~\ref{sec:lucas} или учёт
множественности приложения~\ref{app:lookelsewhere}, которые являются абстрактными
спецификациями, не зависящими от какой-либо конкретной реализации RTL.

\section*{Благодарности}
Laslo Hunhold (Кёльнский университет) предоставил подробную обратную связь по
версии v1.7 данной рукописи, которая сформировала настоящую ревизию: более
короткую двухчастную аннотацию; расширенное введение, обрамляющее компромисс
динамического диапазона и мотивирующее выведенное из \(\varphi\) разбиение до
формулировки правила; абзац мотивации (потеря значимости при большом накоплении
и сравнение с quire/Кулиш), помещённый перед предложением~\ref{prop:lucas};
использование \emph{предложения} вместо эпистемических меток в основном тексте;
определение лестницы GF при \(N \ge 4\) с \(N \in \{2, 3\}\), отмеченными как
вырожденные краевые случаи формулы; и единообразную терминологическую замену
\emph{мантиссы (fraction)} на \emph{мантиссу} повсюду. Ни одна из оставшихся
переоценок, ошибок или открытых гипотез не должна приписываться Hunhold; они
являются ответственностью автора.

% ======================================================================
% REFERENCES
% ======================================================================
\begin{thebibliography}{99}\small

\bibitem{bergman1957} G. Bergman, ``A number system with an irrational base,''
\emph{Mathematics Magazine}, vol.~31, pp.~98--110, 1957.

\bibitem{daubechies2010gre} I. Daubechies, C. S. Gunturk, Y. Wang, and
O. Yilmaz, ``The Golden Ratio Encoder,'' \emph{IEEE Transactions on
Information Theory}, vol.~56, no.~10, pp.~5097--5110, 2010.
DOI 10.1109/TIT.2010.2059750. arXiv:0809.1257.

\bibitem{ward2008} R. Ward, ``On robustness properties of beta encoders and
golden ratio encoders,'' 2008. arXiv:0806.1083.

\bibitem{gustafson2017posit} J. L. Gustafson and I. Yonemoto, ``Beating
floating point at its own game: posit arithmetic,'' \emph{Supercomputing
Frontiers and Innovations}, vol.~4, no.~2, 2017.
DOI 10.14529/JSFI170206.

\bibitem{positstd2022} Posit Working Group, ``Standard for Posit Arithmetic,''
2022. \texttt{https://posithub.org/docs/posit\_standard-2.pdf}.

\bibitem{dedinechin2019} F. de Dinechin, L. Forget, J.-M. Muller, and Y.
Uguen, ``Posits: the good, the bad and the ugly,'' 2019.
\texttt{hal-03195756v3}.

\bibitem{hunhold2024takum} L. Hunhold, ``Takum arithmetic,'' 2024.
arXiv:2404.18603.

\bibitem{hunhold2024int} L. Hunhold, ``Integer representations of takums,''
2024. arXiv:2412.20273.

\bibitem{hunhold2024vhdl} L. Hunhold, ``A VHDL codec for takum arithmetic,''
2024. arXiv:2408.10594.

\bibitem{hunhold2024linear} L. Hunhold and S. Quinlan, ``Bfloat16, posit and
takum number formats in linear-solver workloads,'' 2024. arXiv:2412.20268.

\bibitem{hunhold2025arith} L. Hunhold, ``Hardware evaluation of takum
arithmetic,'' \emph{ARITH 2025 (IEEE 32nd Symposium on Computer Arithmetic)}.
DOI 10.1109/ARITH64983.2025.00019.

\bibitem{mallasen2022percival} D. Mallasen et al., ``PERCIVAL: open-source
posit RISC-V core,'' \emph{IEEE Transactions on Emerging Topics in
Computing}, 2022. DOI 10.1109/TETC.2022.3187199. arXiv:2111.15286.

\bibitem{mallasen2024bigpercival} D. Mallasen et al., ``Big-PERCIVAL:
exploring the native use of 64-bit posit arithmetic,'' \emph{IEEE
Transactions on Computers}, 2024. DOI 10.1109/TC.2024.3377890.
arXiv:2305.06946.

\bibitem{rouhani2023mx} B. D. Rouhani et al., ``Microscaling data formats
for deep learning'' (OCP-MX), 2023. arXiv:2310.10537.

\bibitem{lns_madam2021} J. Zhao et al., ``LNS-Madam: low-precision training
in logarithmic number system,'' 2021. arXiv:2106.13914.

\bibitem{alam2021lns} M. S. Alam, J. Garland, and D. Gregg, ``Low-precision
logarithmic number systems,'' 2021. arXiv:2102.06681.

\bibitem{ahlbach2012} C. Ahlbach, J. Usatine, and N. Pippenger, ``Efficient
algorithms for Zeckendorf arithmetic,'' 2012. arXiv:1207.4497.

\bibitem{alsuhli2023survey} H. F. Alsuhli et al., ``Number systems for deep
neural network architectures: a survey,'' 2023. arXiv:2307.05035.

\bibitem{micikevicius2022fp8} P. Micikevicius et al., ``FP8 formats for deep
learning,'' 2022. arXiv:2209.05433.

\bibitem{luo2024hifloat8} Y. Luo et al., ``Ascend HiFloat8 format for deep
learning,'' 2024. arXiv:2409.16626.

\bibitem{kloewer2020} M. Kloewer, P. D. Duben, and T. N. Palmer, ``Number
formats, error mitigation, and scope for 16-bit arithmetics in weather and
climate modeling,'' \emph{Journal of Advances in Modeling Earth Systems},
vol.~12, 2020. DOI 10.1029/2020MS002246.

\bibitem{bitnet2024} S. Ma et al., ``The era of 1-bit LLMs: BitNet b1.58,''
2024. arXiv:2402.17764.

\bibitem{gustafson2015endoferror} J. L. Gustafson, \emph{The End of Error:
Unum Computing}. CRC Press, 2015. ISBN 9781482239874.

\bibitem{kumar2024scaling} T. Kumar et al., ``Scaling laws for precision,''
2024. arXiv:2411.04330.

\bibitem{bitnet2b4t2025} S. Ma et al., ``BitNet b1.58 2B4T technical
report,'' 2025. arXiv:2504.12285.

\bibitem{fibbinary2025} M. I. Belghazi et al., ``Fibbinary: quantization
of neural networks using Fibonacci representations,'' 2025.
arXiv:2511.01921.

\bibitem{p3109_v091} IEEE P3109 Working Group, ``Standard for binary
floating-point arithmetic for machine learning,'' working draft v0.9.1, 2025.
Reference implementation: \texttt{graphcore-research/gfloat}.

\bibitem{kulisch2013} U. Kulisch, \emph{Computer Arithmetic and Validity:
Theory, Implementation, and Applications}, 2nd ed. De Gruyter, 2013.
ISBN 9783110301731.

\end{thebibliography}

% ======================================================================
% АНГЛОЯЗЫЧНЫЙ БЛОК (требование «Программирование» / Programming and Computer Software)
% ======================================================================
\clearpage
\selectlanguage{english}
\begin{center}
{\large\textbf{GoldenFloat: a family of static-split floating-point formats\\
derived from the golden ratio \(\varphi\), from GF4 to GF256,\\
with an exact integer Lucas identity}}\\[6pt]
\textbf{D.~V.~Vasilev}\\[4pt]
\emph{Independent researcher, Ko Samui, Thailand}\\
e-mail: admin@t27.ai \quad ORCID: 0009-0008-4294-6159
\end{center}

\medskip
\noindent\textbf{Abstract.}
We present a hardware-oriented description of GoldenFloat~(GF) --- a family of
static-split floating-point formats generated by a single closed-form rule ---
together with three concrete artefacts: (i)~an open multi-width RTL generator
covering GF4--GF256 with a continuous-integration differential sweep against a
correctly-rounded reference; (ii)~a Lucas-exact integer accumulator path verified
to 500 digits for \(n = 1,\dots,256\); and (iii)~a GF16 FPGA codec passing a
35/35 testbench at 323~MHz on Artix-7 (Xilinx~XC7A35T). For every full width
\(N \ge 4\) the exponent-field width is \(e = \mathrm{round}((N-1)/\varphi^{2})\)
with fraction width \(f = N-1-e\) and \(\varphi = (1+\sqrt{5})/2\). The rule
reproduces the realised exponent widths of nine formats (9/9) and extends
consistently to further widths. It is positioned alongside posit (Posit
Standard 2022), takum (Hunhold 2024, 2025), OCP-MX (Rouhani et al.\ 2023) and
the IEEE~P3109 multi-width floating-point draft. The breadth/toolchain-coherence
thesis is registered as an open conjecture with a pre-registered falsification
path (ledger FL-002); we make no per-step accuracy or superiority claim. An RTL
correctness erratum (2026-05-31) is reported: the fabricated TTSKY26b dies carry
a defective multiplier portfolio, and the corrected generator is the baseline
for regeneration.

\medskip
\noindent\textbf{Keywords:} numeric formats; floating point; golden ratio;
Lucas numbers; posit; takum; OCP-MX; arithmetic hardware.

\medskip
\noindent\emph{References: see the bibliography list above; the English
``References'' list is produced by the publisher from the same entries.}
\selectlanguage{russian}

\end{document}
